\section{Discussion and open problems}\label{sec:discussion}
The hierarchy separates four ways to ask for certainty. Finite-dimensional perfection has a finite existential witness. Commuting perfection is controlled by compactness. Value one asks for arbitrarily accurate finite witnesses. Spatial perfection combines an unbounded finite choice with an infinite verification. Here the chosen object is the order of a clock phase.

The entropy restriction does not lower the complexity. Each YES instance $G_e$ of our reduction has a finite-entropy witness. There is no computable entropy budget sufficient for every finite-entropy YES game: such a budget $B(G)$ would make $\perf_{\finS}$ a $\Pi^0_1$ set by the uniform budget predicate in Section~\ref{sec:upper}, contradicting its $\Sigma^0_1$-hardness from perfect completeness of $\mathsf{MIP}^*=\mathsf{RE}$. The analogous impossibility for a sufficient largest-Schmidt stratum is Theorem~C. The fixed-$\kappa$ compactness result does not give a uniform modulus over games. It also concerns zero constraints and extremal scalar eigen-equations, not a full prescribed correlation table.

\subsection{Prescribed tables and escaping mass}
For a concrete obstruction to the latter extension, put
\[
 D_{{\rm esc},n}=\operatorname{diag}(1,\underbrace{n^{-1/2},\ldots,n^{-1/2}}_{n\text{ entries}},0,\ldots).
\]
Then $\Tr D_{{\rm esc},n}^2=2$ and $D_{{\rm esc},n}\to|0\rangle\langle0|$ in operator norm. Alice measures $|0\rangle\langle0|$ and its complement; Bob has one trivial outcome. Every normalized vectorization of $D_{{\rm esc},n}$ gives Alice probabilities $(1/2,1/2)$. The normalized limit gives $(1,0)$. The largest-amplitude normalization preserves zero equations while changing nonzero normalized probabilities.

\paragraph{Unproved remarks on prescribed-table bounds.}
The following upper-bound remarks are not proved in this paper and are not
used in any theorem. For a prescribed table with rational or real-algebraic entries, direct coding of countable operators and a normal vector gives a $\Sigma^1_1$ upper bound. A supplied computable Schmidt-tail modulus makes the normalized state family effectively compact and gives a $\Pi^0_1$ upper bound. An existential choice from a uniformly total computable family of such moduli gives $\Sigma^0_2$. Finite entropy also gives $\Sigma^0_2$, by \eqref{eq:entropytail} with an existential integer budget and rational slack. If an arbitrary program for a computable modulus is existentially supplied, totality and the tail conditions give a $\Sigma^0_3$ upper bound. These proposed bounds concern the respective restricted realization classes; establishing their effective coding details is outside the present proofs. We do not prove that every algebraic spatial table has a computable-tail realization, or give a completeness classification for unrestricted table membership.

\subsection{The Fu--Miller--Slofstra remark}
After Theorem~1.1, Fu, Miller and Slofstra state that their fixed-scenario membership problem is in $\mathsf{RE}$ for $t=q$ or $t=qs$ \cite{FMS}. We found no proof of the general $qs$ upper assertion in that discussion. The following consequence needs a uniform interpretation over finite scenarios.

Suppose membership of real-algebraic tables in $C_{qs}$ has a uniform RE procedure. If every spatially perfect game in our reduction had an algebraic perfect table, enumerate all algebraic tables of its scenario, test the losing zeros exactly, and dovetail the membership procedures. This would enumerate precisely the YES inputs of a $\Sigma^0_2$-complete family, a contradiction. Under that interpretation, either the RE remark fails for $qs$, or some spatially perfect games in this family have no perfect correlation table with all entries algebraic. The conclusion is existential over YES games; it does not say that every YES game has only nonalgebraic perfect tables.

The literal remark is per fixed scenario. The scenario of $G_e$ grows with $e$, which is hardwired into the CA alphabet, so separate nonuniform RE procedures for each scenario give no contradiction. Tension arises only if there is one procedure uniform across scenarios, together with algebraic perfect tables for every YES game of the reduction. This is an open discrepancy, not a refutation. Resolving the uniformity and algebraicity questions is a natural next problem.

\subsection{The von Neumann type ladder}\label{sec:typeladder}
Version 1.0.0 of this paper asked how perfect play behaves when the ambient local algebras have type II or type III. Three papers by the authors answer them. The results below are proved there, not here. In each, a strategy is perfect normal play on an ambient von Neumann algebra of the stated type, with the commuting partner specified in the cited paper. The spatial strategies of Theorem~A are the type-I case.

A single finite game already separates type III from semifinite hosts. The game $G_{\rm III}$ of \cite[Theorem~1.1]{ReadType} has $45$ Alice questions, $46$ Bob questions and $243$ uniformly sampled question pairs. It has a perfect normal strategy on the injective type $\mathrm{III}_{1/2}$ factor, and no perfect strategy has a normal marginal on a semifinite ambient algebra of either party. In particular it has perfect commuting play but no perfect spatial strategy. For factors with separable predual, hosting $G_{\rm III}$ is equivalent to containing an expected copy of the injective type $\mathrm{III}_1$ factor. Injective hosts are excluded by a different game. The game of \cite[Theorem~1.1]{Hyperfinite} has a perfect correlation in the closure of the finite-dimensional correlations, realized on a matrix amplification of a type $\mathrm{III}_{1/2}$ factor. No perfect strategy admits a normal injective host for either player. Normal hyperfinite strategies therefore approach perfect play for this game but never attain it.

The arithmetical classification is \cite[Theorem~A]{TypeLadder}. For finite rational games, perfect normal play in a type $\mathrm{II}_1$ algebra, in a type $\mathrm{II}_\infty$ algebra, or in some semifinite algebra is $\Sigma^0_2$-complete; the upper bound is Theorem~2.2 of \cite{TypeLadder}, and the hardness reduction is Theorem~C there. For each fixed $\lambda\in[0,1]$, the type $\mathrm{III}_\lambda$ language equals commuting perfection \cite[Theorem~7.2]{TypeLadder} and is therefore $\Pi^0_1$-complete. Commuting perfection without semifinite perfection is $\Pi^0_2$-complete, with hardness on a computable family of commuting-perfect games of quantum value one. The type-I language, whose classification is Theorem~A of the present paper, and the type $\mathrm{II}_1$ language are incomparable, and both are proper subsets of the type $\mathrm{II}_\infty$ language. Theorem~9.1 and Corollary~9.2 of \cite{TypeLadder} show that two directional differences are $D_2(\Sigma^0_2)$-complete: perfect play in type I but not in type $\mathrm{II}_1$, and perfect spatial play attainable only with infinite entropy. Completeness of semifinite-but-not-type-I perfection, and of its finite-entropy variant, remains open there; both are $\Sigma^0_2$-hard and $\Pi^0_2$-hard. These are existence languages, in which the ambient algebra and its normal state are chosen together, with the conventions fixed in \cite[Section~1]{TypeLadder}. Separations involving spatial play there use the reduction of Theorem~\ref{thm:reduction}.

\subsection{Other questions}\label{sec:otherquestions}
Can the game compiler be made substantially smaller? Can the period reader be realized by a simpler predicate? What are the optimal effective bounds when one supplies an entropy budget or a Schmidt-tail modulus? The finite-entropy and arbitrary-spatial decision sets have the same completeness degree, although the sets themselves differ: the entropy-game construction in \cite{Entropy} has spatial perfection and no finite-entropy perfect strategy. Their difference, the set of games with perfect spatial play but no perfect spatial strategy of finite entropy, is $D_2(\Sigma^0_2)$-complete \cite[Corollary~9.2]{TypeLadder}; in particular it belongs to neither $\Sigma^0_2$ nor $\Pi^0_2$. A classification of prescribed-table membership would require additional control of the normalization that escapes in the example above.
